\chapter{Discussion and Conclusion}

\section{Summary of Findings}

The experimental results presented in Chapter~2 yield five key findings:

\textbf{1. Decision-making remains challenging.} Across GPT-5-mini, DeepSeek-v3.2, and Qwen-Turbo, average decision-making performance remains around 0.42--0.45, with substantial variation across subtasks. The Discharge Decision is the easiest subtask (roughly 0.58--0.63), while Next Action Prediction remains much lower (roughly 0.28--0.34). This gap shows that binary discharge judgment should not be treated as a proxy for full clinical decision-making ability.

\textbf{2. Memory granularity matters modestly.} Structured Event Memory slightly outperforms Note Memory in most conditions, but the effect is small relative to the overall performance ceiling. For proactive clinical planning, immediate context dominates.

\textbf{3. Memory augmentation does not improve decision-making.} This is the central finding. Neither RAG nor Mem0 provides significant improvement over the naive long-context baseline. In some conditions, augmentation slightly degrades performance.

\textbf{4. Immediate context dominates historical context.} The $n=0$ baseline ($\mathcal{M}_C^{(j,T)}$ only, no history) achieves accuracy comparable to full-history conditions. Adding historical visits does not monotonically improve decisions.

\textbf{5. Chain-of-thought can degrade performance.} For DeepSeek-v3.2 under Event Memory, reasoning mode reduces average decision-making performance from 0.439 to 0.419 and Discharge Decision accuracy from 0.607 to 0.593, likely because explicit reasoning induces conservatism that penalizes decisive action under multiple-choice evaluation.

\section{The Retrieval-Reasoning Gap}

These findings collectively reveal a \textbf{retrieval-reasoning gap}: a disconnect between LLMs' ability to \emph{locate} relevant historical information and their ability to \emph{integrate} that information into proactive clinical planning. This explains why RAG and Mem0 fail: these architectures optimize retrieval quality, but for clinical decision-making, retrieval is not the bottleneck. The model can access relevant history (either in full context or via retrieval), but cannot effectively use it to inform decisions.

We hypothesize three contributing mechanisms. First, clinical reasoning requires causal, not associative, inference---determining whether a prior treatment was effective requires understanding why it was given and how the patient responded, not just retrieving that it occurred. Second, the signal-to-noise ratio in longitudinal EHR data is low: a patient's 900 events contain only a fraction relevant to any single decision, overwhelming attention mechanisms. Third, current LLM training data (web text, code, books) lacks the temporal-causal structure of clinical trajectories, creating a distribution shift.

\section{Implications}

\subsection{For Memory Architecture Design}

Current memory architectures (RAG, Mem0, MemGPT) operate on the assumption that better retrieval yields better performance. Our findings challenge this assumption for clinical decision-making. Future architectures should emphasize \emph{reasoning over retrieved information}: temporal attention mechanisms that encode causal relevance, graph-based representations that capture clinical relationships, and structured reasoning layers that bridge retrieval and decision.

\subsection{For Benchmark Design}

If decision-making performance is dominated by immediate context, benchmarks that primarily test retrieval may overestimate clinical readiness. We advocate for adversarial task design---constructing decision instances where history \emph{changes} the correct answer---to prevent models from succeeding through immediate-context heuristics alone. Additionally, distinguishing the subset of decisions where history genuinely matters from those driven by current presentation is critical for targeted evaluation.

\subsection{For Clinical AI Deployment}

The finding that immediate context dominates is not necessarily a weakness---in many clinical scenarios, the current presentation is indeed the most important determinant. The challenge lies in identifying decisions where ignoring history constitutes clinical error (e.g., prescribing a medication that previously caused an adverse reaction) and ensuring the agent has the architecture to recognize and act on such cases.

\section{Limitations}

\textbf{Data Scope.} LongMedBench is constructed from MIMIC-IV, a single-institution database. Clinical practice patterns and documentation conventions vary across institutions, and findings may not generalize without cross-institutional validation.

\textbf{Cohort Bias.} Several design decisions in our data pipeline shape the cohort composition. At the event level, we retain only laboratory results flagged as abnormal, following Hager et al.\ \cite{hager2024evaluation}, who demonstrated that presenting models with only abnormal lab values improved diagnostic performance in MIMIC-CDM by reducing noise from clinically uninformative normal results. This filtering makes the evaluation more clinically realistic---like human clinicians, the agent reasons from salient pathological signals---but may underrepresent the diagnostic challenge of distinguishing normal variation from true abnormality. At the patient level, we require each hospitalization to have a matched discharge note from the \texttt{note} module and select patients with $\ge 15$ hospitalizations. We chose this threshold after examining the trade-off between trajectory length and cohort size: lowering to $\ge 10$ visits yields trajectories too short to meaningfully evaluate cross-visit reasoning; raising to $\ge 20$ visits produces trajectories of over 20 steps but retains only approximately 100 patients, limiting statistical power. At $\ge 15$ visits, the resulting cohort of 355 patients compares favorably to sample sizes in other EHR-based agent benchmarks---MedAlign enrolled 276 patients \cite{2024_medalign} and DiReCT enrolled 511 patients \cite{wang2024direct}. Nevertheless, this threshold selects for high-acuity, chronic, multi-morbid patients---the population for whom longitudinal reasoning matters most---and findings may not generalize to lower-acuity settings or the broader patient base.

\textbf{Task Simplification.} The benchmark decomposes clinical decision-making into three multiple-choice subtasks: Next Action Prediction predicts the next event action type, Argument Prediction predicts the event argument, and Discharge Decision predicts six-hour discharge readiness. This decomposition improves diagnostic clarity but still simplifies real clinical decisions, which involve dosing, contraindication checking, patient preferences, resource availability, and guideline compliance not fully captured by our formulation.

\textbf{Question Discriminability.} Not all evaluation questions are equally informative. An analysis of model agreement reveals that the Discharge Decision subtask has extremely few discriminative instances: only 8 out of 448 sampled questions (~1.8\%) produce disagreement across model configurations; the vast majority yield unanimous scores of 1.0 because patients are simply too early or too late for discharge. This inflates apparent Discharge Decision accuracy and limits its diagnostic value for comparing architectures. The Next Action Prediction subtask, by contrast, has approximately 70\% discriminative questions, making it the most informative component of the decision-making benchmark.

\textbf{Evaluation Rigor.} Classification accuracy for Next Action Prediction and Discharge Decision and top-$k$ accuracy for Argument Prediction may penalize clinically reasonable alternatives. The time-decay mechanism partially addresses timing ambiguity for Next Action Prediction and Argument Prediction, but no automatic metric replaces clinical expert judgment. We recommend human clinician evaluation as a future validation step.

\textbf{Model Coverage.} Our evaluation covers three general-purpose LLMs. Models specifically trained for clinical tasks (e.g., EHR-R1, Med-PaLM) may exhibit different patterns not captured here.

\section{Future Work}

\textbf{Task Refinement.} Adversarial decision instances where history changes the correct answer would directly test temporal integration. Counterfactual evaluation---altering history and measuring decision sensitivity---would isolate causal effects.

\textbf{Memory Architecture Development.} Temporal graph memory, selective history summarization (compressing irrelevant while preserving salient information), and memory-aware fine-tuning on synthetic longitudinal trajectories are promising directions.

\textbf{Human Baseline.} A clinician evaluation study contextualizing both average decision-making performance (approximately 0.42--0.45) and subtask-specific performance would validate benchmark difficulty and enable calibrated assessment of model progress.

\textbf{Cross-Institutional Validation.} Extending the pipeline to eICU, HiRID \cite{2021_hirid_icu}, or STARR would test generalizability across institutions and patient populations.

\textbf{Multimodal Extension.} Incorporating medical imaging and physiological waveforms would increase clinical fidelity as vision-language and multimodal LLMs mature.

\textbf{Outcome Correlation.} Linking agent decisions to simulated patient outcomes (mortality, readmission, length of stay) would provide the ultimate validation of decision-making quality.

\section{Conclusion}

Clinical practice is inherently longitudinal---patients accumulate histories over years, and clinicians must integrate evidence across this temporal expanse. As LLMs move toward clinical deployment, their capacity for longitudinal reasoning becomes a critical benchmark.

This thesis presented the decision-making component of LongMedBench, investigating how contemporary LLMs use historical clinical information to inform proactive planning through a three-tier memory architecture and systematic ablation experiments. The central finding is that memory augmentation---the dominant paradigm for extending LLM capabilities---does not significantly improve longitudinal clinical decision-making. Decision accuracy depends primarily on immediate context, and historical information, however organized, contributes marginal additional value.

This finding clarifies the nature of the challenge. The bottleneck is not retrieval---finding the right facts---but integration---understanding how those facts collectively inform decisions. Closing this retrieval-reasoning gap requires architectures that go beyond retrieval to support structured reasoning over temporal clinical evidence. The path to clinically capable AI agents is longer than single-encounter benchmarks suggest, but by identifying where current approaches fall short, LongMedBench provides a platform for measuring and ultimately closing this gap.

\subsection{Relationship to the Broader LongMedBench Project}

While this thesis focuses on the Long-Horizon Decision-Making task, placing our findings in the context of the broader LongMedBench project reveals complementary insights. LongMedBench also includes a Factual QA task, which tests factual retrieval from long clinical histories, and a Temporal Reasoning task, which tests event and visit ordering. These tasks help separate retrieval, temporal reconstruction, and decision-making.

The contrast between these task suites sharpens our central finding: memory augmentation can help retrieval-oriented tasks, but it does not reliably improve decision-making across Next Action Prediction, Argument Prediction, and Discharge Decision subtasks. This dissociation confirms that the retrieval-reasoning gap is task-specific---it manifests when the agent must do more than locate facts, when it must convert longitudinal evidence into an actionable decision. This insight would not be visible from any single task suite; it emerges only from the systematic, multi-dimensional evaluation that LongMedBench provides.

\subsection{Practical Considerations for Reproducibility}

The MIMIC-IV data itself requires credentialed access through PhysioNet. To facilitate reproducibility and future extensions, we will release the data processing pipeline, task generation logic, and evaluation harness, if LongMedBench is admitted by MICCAI 2026. 

API-based model evaluation used default settings and fixed random seeds where applicable. Full prompt templates, including system prompts, memory formatting conventions, and decision queries, are documented in the released codebase to ensure exact replication.
